\documentclass[11pt,a4paper]{article}
\usepackage[T1]{fontenc}
\usepackage[utf8]{inputenc}
\usepackage{lmodern}
\usepackage{amsmath,amssymb,amsthm,mathtools,mathrsfs}
\usepackage{graphicx,float,xcolor,multicol}
\usepackage[shortlabels]{enumitem}
\usepackage[margin=27mm]{geometry}
\usepackage{microtype,needspace,placeins}
\usepackage{tikz,tikz-cd}
\usetikzlibrary{arrows.meta,calc,positioning,patterns,decorations.pathreplacing}
\usepackage[colorlinks=true,linkcolor=teal,urlcolor=teal,citecolor=teal]{hyperref}
\setlength{\parindent}{0pt}
\setlength{\parskip}{.6em}
\setcounter{secnumdepth}{0}
\allowdisplaybreaks
\raggedbottom
\providecommand{\cross}{\times}
\DeclareUnicodeCharacter{2020}{\ensuremath{\dagger}}
\DeclareMathOperator{\supp}{supp}
\DeclareMathOperator*{\esssup}{ess\,sup}
\providecommand{\chapter}{\section}
\newenvironment{tool}[2]{\subsection*{Tool #1 --- #2}}{}
\newcommand{\ToolRef}[1]{Tool #1}
\newcommand{\FigureTag}[2]{\textit{Figure #1}}
\newcommand{\Result}[1]{\subsection*{#1}}
\newcommand{\Application}[1]{\subsubsection*{Application\if\relax\detokenize{#1}\relax\else\space--- #1\fi}}
\newcommand{\ProofHeading}{\par\textit{Proof.}\space}
\newcommand{\ProofOf}[1]{\par\textit{Proof of #1.}\space}
\newcommand{\RemarkHeading}{\par\textbf{Remark.}\space}
\newcommand{\NoteHeading}{\par\textbf{Note.}\space}
\newcommand{\ExerciseHeading}{\par\textbf{Exercise.}\space}

\graphicspath{{figures/k-theory/}}
\begin{document}
\section*{Classification of Vector Bundles}
% BLOG-CONTENT-BEGIN
Of abundant importance to any theory is a robust classification theorem. With that interest in mind let us prove the following theorem pointing in that direction for vector bundles.
\begin{theorem}[Homotopy theorem]
    Given a vector bundle $\eta=(E,B,V,p)$ and homotopic maps $f_0,f_1:A\to B$, where $A$ is paracompact Hausdorff, we have:
    $$f_0^*\eta \cong f_1^*\eta$$
\end{theorem}

\begin{proof}
    We shall do it in three steps, each important in itself.

    \subsection*{Bundle maps and sections} $\text{Hom}(E, E') \longleftrightarrow \Gamma(\text{HOM}(E, E'))$.

    Here the specific description of the HOM bundle is useful. Define,
    $$\rho: \text{Aut}(V)\times \text{Aut}(V')\rightarrow \text{Aut}(\text{Hom}(V, V'))$$
    \begin{equation}\tag{1}
        \rho(T, S)(A) := SAT^{-1}
    \end{equation}

The functorial construction yields the $\rho$ above, that is,
$$\rho(E, E') = \text{HOM}(E, E')$$
Now let $f\in \text{Hom}(E, E')$ and let $\{\varphi^1_{\alpha}\}_{\alpha}$ and $\{\varphi^2_{\alpha}\}_{\alpha}$ be local trivializations of $E$ and $E'$ (after passing to a common refinement in $B$). We have,
$$\big((\varphi^2_{\alpha})^{-1}\circ f
\circ \varphi^1_{\alpha}\big)(x, h) = (x, f_{\alpha}(x)h)$$
$$\widetilde{\varphi}^2_{\beta\alpha}\circ f_{\alpha} \equiv f_{\beta}\circ \widetilde{\varphi}^1_{\beta\alpha}$$
Clearly, $f_{\alpha}$ is a local section of HOM$(E, E')$ because if $\{\varphi_{\beta\alpha}\}_{\alpha, \beta}$ are transition functions of it, then by $(1)$,
$$f_{\beta} = \varphi_{\beta\alpha}\circ f_{\alpha}$$
which implies that $\{f_{\alpha}\}_{\alpha}$ can be glued together on $B$ to yield a global section.

\subsection*{Extending a bundle isomorphism} Let $B$ be paracompact Hausdorff, let $B'\subseteq B$ be closed, and let $E\big|_{B'} \rightarrow E'\big|_{B'}$ be an isomorphism. Then there exists $U\subseteq_{\text{open}} B$, $B'\subseteq U$ such that the isomorphism can be extended to $E\big|_{U} \rightarrow E'\big|_{U}$.

We shall now make use of the following:

\underline{FACT}: A section of a vector bundle over a closed subspace of a paracompact Hausdorff space extends to a section over an open neighbourhood of that subspace.

By the bundle-maps-and-sections correspondence, $f\in \Gamma\big(\text{HOM}(E\big|_{B'}, E'\big|_{B'}\big)$, and by the fact above, $f$ can be extended to $\widetilde{f}$ over an open neighbourhood of $B'$. Also, the set of points $x\in B$ for which $\widetilde f(x)$ is an isomorphism is open in $B$.

\subsection*{Applying the extension to a homotopy} Let $F: A\times I\rightarrow B$ be the homotopy, let $F_t:= F(\cdot, t)$, and let $\pi:A\times I\rightarrow A$ be the projection.
$$\implies F_t\circ \pi : A\times I \rightarrow B$$
Consider the bundles $F^*E$ and $(F_t\circ \pi)^*E$, note that, $F_t\circ \pi = F\big|_{A\times \{t\}}$. Applying the extension result to $F^*E$ and $(F_t\circ \pi)^*E$ restricted to $A\times \{t\} \subseteq_{\text{closed}} A\times I$, we see that all bundles $F^*E\big|_{A\times \{\widetilde{t\}}}$, for $\widetilde{t} \in (t-\delta, t+\delta)$ are isomorphic for some $\delta$. As $[0,1]$ is connected, we get that:
$$f_0^*\eta \cong f_1^*\eta$$
\end{proof}


\begin{corollary}
    If $f:A\rightarrow B$ is a homotopy equivalence, then, $f^*: \text{Vect}^n(B)\rightarrow \text{Vect}^n(A)$ is a bijection. ((i.e.) vector bundles are homotopy invariant).
\end{corollary}
\begin{proof}
    Let $g:B\rightarrow A$ be such that $f\circ g \sim \text{id}_B$ and $g\circ f\sim \text{id}_A$, then $g^*$ is the required inverse for $f^*$, as $f^*g^* \cong \text{id}^*_A$ and $g^*f^* \cong \text{id}^*_B$
\end{proof}

\underline{\textbf{Classification of Vector Bundles} (over compact Hausdorff spaces)}

Let $\tau_c$ be the category of compact Hausdorff spaces up to homotopy. Note that $\text{Vect}_{\mathbb{K}}^n:\tau_c^{\text{op}}\to\text{Sets}$ is a presheaf on $\tau_c$, and hence asking for a good classification theorem is akin to asking if $\text{Vect}_{\mathbb{K}}^n$ is representable. The answer is yes, by the infinite Grassmannian $\text{Gr}(n, \mathbb{K}^{\infty})$. More specifically, one can obtain all vector bundles over $X$ by pulling back via maps $X\rightarrow \text{Gr}(n , \mathbb{K}^{\infty})$ the tautological bundle $\gamma_n$ over $\text{Gr}(n , \mathbb{K}^{\infty})$. For apparent reasons, $\text{Gr}(n , \mathbb{K}^{\infty})$ is called the classifying space of the structure group $GL_n(\mathbb{K})$.
\begin{theorem}[Classification of vector bundles]
    The map $\phi: [X, \text{Gr}(n , \mathbb{K}^{\infty})] \rightarrow \text{Vect}_{\mathbb{K}}^n(X)$ is a bijective correspondence for compact Hausdorff spaces $X$, where $[X, Y]$ denotes the set of maps $X\rightarrow Y$ up to homotopy, and $\phi([f]):= [f^*\gamma_n]$.
\end{theorem}

\begin{remark}
A more general result is true: Given a group $G$, let $\mathcal{P}_G(X)$ denote the set of all principal $G$-bundles, where $X$ is compact Hausdorff, then there exists a $G$-bundle $EG\rightarrow BG$, with a contractible total space, such that the following is a bijective correspondence:
$$[X, BG] \longleftrightarrow \mathcal{P}_G(X)$$
$EG$ is called the \textbf{universal bundle} and $BG$ the \textbf{classifying space} for the structure group $G$. This generalizes the theorem above for $G= GL_n{(\mathbb{K})}$, because of the correspondence between $GL_n{(\mathbb{K})}$-bundles and vector bundles of dimension $n$ established before, and the fact that $BGL_n{(\mathbb{K})} \simeq \text{Gr}(n, \mathbb{K}^{\infty})$.
\end{remark}

\subsection*{Scope of the classification theorem}
Note that the homotopy theorem ensures that the pullback of bundles via maps in $[X, \text{Gr}(n, \mathbb{K}^{\infty})]$ is well-defined. The hypothesis on $X$ can be weakened to allow paracompact Hausdorff spaces and numerable bundles. Under these hypotheses, the following is a bijective correspondence:
$$[X, \text{Gr}(n, \mathbb{K}^{\infty})] \longrightarrow \text{Vect}_{\mathbb{K}}^n(X)$$
$$[f] \rightarrow [f^*\gamma_n]$$

\begin{proof}
    To establish injectivity we need the following lemma:
    \begin{lemma}
        Let $j_{\text{ev}}$, $j_{\text{odd}}$ $: \mathbb{K}^{\infty} \rightarrow \mathbb{K}^{\infty}$ be defined as follows:
        $$j_{\text{ev}}(x_1, x_2, \ldots) := (0, x_1, 0, x_2, 0, \ldots)$$
        $$j_{\text{odd}}(x_1, x_2, \ldots) := (x_1, 0, x_2, 0, \ldots)$$
        Then $j_{\text{ev}} \sim j_{\text{odd}} \sim \text{id}_{\mathbb{K}^{\infty}}$
    \end{lemma}
    
    \begin{proof}
        We prove it for $j_{\text{ev}}$, and for $j_{\text{odd}}$ the proof is analogous. Define,
        $$H: \mathbb{K}^{\infty} \times [0,1] \rightarrow \mathbb{K}^{\infty}$$
        $$H(x,t):= tj_{\text{ev}}(x) + (1-t)x$$
        Clearly, $H$ is the required homotopy. For every $t\in[0,1]$, the equality $H(x,t)=0$ implies $x=(0,0,\ldots)$. So, the homotopy is via linear embeddings.
    \end{proof}
    
    
    Let $f$, $g$ $: X\rightarrow \text{Gr}(n, \mathbb{K}^{\infty})$ be maps, and suppose $\alpha: f^*\gamma_n \xrightarrow{\cong} g^*\gamma_n$ is an isomorphism. By replacing $f$ and $g$ with $j_{\text{ev}}\circ f$ and $j_{\text{odd}}\circ g$, we may assume their images lie in orthogonal summands of $\mathbb K^\infty$. The bundle isomorphism $\alpha:f^*\gamma_n\to g^*\gamma_n$ restricts to an isomorphism $\alpha_x:f(x)\to g(x)$ on every fibre. Now define,
$$\widetilde{H}: X\times [0,1] \rightarrow Gr(n, \mathbb{K}^{\infty})$$
$$\widetilde{H}(x, t):=\big\{(1-t)v+t\alpha_x(v)\mid v\in f(x)\big\}.$$
Clearly, $\widetilde{H}$ is continuous, $\widetilde{H}(x, 0) \equiv f(x)$, and $\widetilde{H}(x, 1) \equiv g(x)$.

Now let's establish surjectivity. Let $E\xrightarrow{p} X$ be a vector bundle of rank $n$. $E$ can be embedded in $X\times \mathbb{K}^N \subset X\times \mathbb{K}^{\infty}$. Let that embedding be $j$. Define,
$$f:X \rightarrow Gr(n, \mathbb{K}^{\infty})$$
$$f(x) := j(E_x) \in Gr(n, \mathbb{K}^{\infty})$$
where $j(E_x) \subseteq x\times \mathbb{K}^N \subseteq x\times \mathbb{K}^{\infty} \cong \mathbb{K}^{\infty}$. Using the local trivialization on $E\rightarrow X$, one can easily prove that $f$ is continuous on each trivialising neighbourhood, and hence continuous. To establish $f^*\gamma_n \cong E$ we need to look at the associated bundle differently. Set,
$$\gamma_n=G_n:= \{(V, w)  |  V \in Gr(n, \mathbb{K}^{\infty}), w\in V\}.$$
Clearly, $f^*G_n \cong E$ via the map:
$$(x,w)\in f^*G_n \rightarrow j_x^{-1}(w)\in E_x$$
with inverse, 
$$v\in E \rightarrow \big(p(v), j(v)\big) \in f^*G_n$$
which is continuous component-wise. $G_n$ is often referred to as the \textbf{tautological bundle}. 
\end{proof}

\begin{remark}
If $p:E\rightarrow B$ is a vector bundle, then one has the fibre preserving sequence $\mathbb{K}^n \rightarrow E\rightarrow B$. This gives rise to a long exact sequence of homotopy groups:
$$\cdots\rightarrow\pi_k(\mathbb{K}^n) \rightarrow \pi_k(E) \rightarrow \pi_k(B) \rightarrow \pi_{k-1}(\mathbb{K}^{n})\rightarrow \cdots$$
As $\mathbb{K}^n$ is contractible, we get the following:
$$\pi_k(E) \cong \pi_k(B)$$
In other words, $E$ is weakly homotopy equivalent to $B$. The above is in fact true for any fibre bundle with contractible fibre, but for the case of vector bundles one can show that the weak equivalence is actually a homotopy equivalence, where the zero section is the prescribed inverse of $p$. This, together with the corollary, to the homotopy theorem and the classification theorem tells us that:
\end{remark}

\begin{center}
    
$\text{Vect}_{\mathbb{K}}^n(B)$ encodes homotopy-theoretic information about $B$ in linear terms.

\end{center}
% BLOG-CONTENT-END
\end{document}
